Este capítulo descreve o projeto desenvolvido: a comparação entre o PostgreSQL com JSONB e o MongoDB em um catálogo de dados parcialmente variáveis. A Seção~\ref{sec:descricao-problema} apresenta o problema e a metodologia, a Seção~\ref{sec:atividades} detalha as atividades realizadas, a Seção~\ref{sec:resultados} analisa os resultados e a Seção~\ref{sec:dificuldades} discute as dificuldades e limitações.

\section{Descrição do Problema e/ou Projeto}
\label{sec:descricao-problema}

Sistemas de catálogo de restaurantes armazenam itens, opções e promoções cujas
estruturas variam de um item para outro. Esses dados podem ser representados em um banco relacional com colunas
JSONB, como o PostgreSQL, ou em um banco orientado a documentos, como o MongoDB.
A pergunta que este trabalho busca responder é: em um sistema de catálogo com dados parcialmente variáveis, em quais cenários o PostgreSQL, com o recurso JSONB, e o MongoDB apresentam melhor desempenho e escalabilidade? 


Para isso, foi implementado nos dois bancos um recorte controlado de um sistema
de catálogo, com dados sintéticos e consultas que foram medidas em
condições idênticas. A Figura~\ref{fig:metodologia} apresenta as etapas da
metodologia, agrupadas em três fases.

\begin{figure}[H]
    \centering
    \caption{Metodologia Geral do Desenvolvimento.}
    \includegraphics[width=0.9\textwidth]{metodologia.png}
    \label{fig:metodologia}
    \legend{Fonte: Elaborado pela autora.}
\end{figure}

Na \textbf{preparação do ambiente e dos dados} (etapas 1 a 4), os dois bancos são
configurados em contêineres Docker, o modelo de dados é definido em cada um e os
mesmos dados sintéticos são gerados e inseridos. Na \textbf{execução do
experimento} (etapas 5 a 7), um conjunto de consultas e operações de escrita é
implementado no PostgreSQL e no MongoDB, avaliado com e sem índices e medido quanto ao
tempo, ao consumo de recursos e ao espaço em disco. Por fim, na
\textbf{escalabilidade e análise} (etapas 8 e 9), as operações são repetidas
variando o volume de dados, a seletividade e o tamanho do lote de escrita, e os
resultados são comparados. O detalhamento de cada etapa é apresentado na seção
seguinte.


\section{Descrição das Atividades Realizadas}
\label{sec:atividades}

\subsection{Ambiente e Ferramentas Utilizadas}

Todo o experimento foi executado localmente, em uma única máquina, com os dois bancos de dados em
contêineres Docker. A implementação está disponível em duas formas: um notebook Jupyter, que apresenta
cada etapa de forma explicativa, e um script Python (\texttt{analises.py}), que executa os experimentos
de ponta a ponta e grava os resultados em arquivos JSON. O código completo, incluindo a implementação do experimento,
o arquivo de configuração dos contêineres, o arquivo de geração de gráficos e lista de dependências, está disponível em um
repositório público no GitHub\footnote{\url{https://github.com/RauanySecci/TCC_Rauany}}.

A máquina utilizada possui processador Intel Core i5-9300H (2,40~GHz), 8~GB de memória RAM e sistema operacional
Ubuntu 22.04.5 LTS. Os contêineres foram executados com Docker Engine, versão 29.1.3.

Os bancos foram definidos a partir de imagens com versão fixada: PostgreSQL 14.18 e MongoDB 8.0.29.
Cada contêiner foi limitado a 2 CPUs e 3~GB de memória, o mesmo limite para os dois, garantindo
condições equivalentes de execução. Os dados ficam em
volumes persistentes, e as portas padrão de cada banco (5432 e 27017) são expostas ao código.

A implementação foi desenvolvida no Visual Studio Code\footnote{\url{https://code.visualstudio.com}}
com Python\footnote{\url{https://www.python.org}} 3.10.12. O notebook e o script concentram as etapas do
experimento: criação do esquema, geração e inserção dos dados, execução das consultas, criação dos
índices, escalabilidade, medição de desempenho e gravação dos resultados. Antes de cada repetição de uma medição, os
contêineres são reiniciados e o cache do sistema operacional é descartado, sem aquecimento, de modo que
cada execução parte do cache vazio. As credenciais e os parâmetros de conexão não ficam no
código: são lidos de arquivos \texttt{.env}.

A primeira etapa do script (e a primeira célula do notebook) realiza a configuração do ambiente, importa as bibliotecas, carrega as variáveis dos arquivos \texttt{.env}, cria a conexão com o
PostgreSQL por meio do SQLAlchemy\footnote{\url{https://www.sqlalchemy.org}} e a conexão com o MongoDB por meio do PyMongo\footnote{\url{https://pymongo.readthedocs.io}}.

As versões de todas as dependências estão fixadas no arquivo \texttt{requirements.txt}, que permite recriar o ambiente
Python e inclui também as dependências indiretas do
Jupyter. As principais bibliotecas utilizadas estão na Tabela~\ref{tab:versoes}, no Apêndice~\ref{chapter:apendice-versoes}.

\subsection{Modelagem dos Dados e Schema}

Para representar o catálogo, o experimento utiliza duas entidades: restaurantes e itens. Essa escolha é um recorte do problema real, já que
em um catálogo de produção há outras entidades e atributos, mas o objetivo aqui é avaliar o comportamento dos bancos sobre dados parcialmente 
variáveis, e não reproduzir o domínio completo. A tabela de itens é a principal do experimento, é nela que se concentram os atributos variáveis, o maior volume de registros 
e a maioria das consultas. A tabela de restaurantes existe para viabilizar a consulta que envolve junção, relacionando os itens ao restaurante que pertencem.

A Figura~\ref{fig:modelo-relacional} apresenta o diagrama UML do modelo adotado no PostgreSQL, onde os dados fixos do item ocupam colunas tradicionais, e os dados que variam ocupam 
colunas JSONB.

\begin{figure}[H]
    \centering
    \caption{Modelo de dados no PostgreSQL.}
    \includegraphics[width=0.8\textwidth]{modelo-relacional.png}
    \label{fig:modelo-relacional}
    \legend{Fonte: Elaborado pela autora.}
\end{figure}

A Figura~\ref{fig:modelo-documentos} apresenta o diagrama UML do modelo equivalente no MongoDB, onde restaurante e item são documentos. Item contém os objetos de preço 
e de grupos de opções e recebe os atributos variáveis como campos próprios (exemplificado por \texttt{atributosEspecificosCategoria}).

\begin{figure}[H]
    \centering
    \caption{Modelo de documentos no MongoDB.}
    \includegraphics[width=0.9\textwidth]{modelo-documentos.png}
    \label{fig:modelo-documentos}
    \legend{Fonte: Elaborado pela autora.}
\end{figure}

A Figura~\ref{fig:hierarquia-item} mostra a hierarquia conceitual dos itens, onde 
todo item pertence a uma das quatro categorias (pizza, bebida, sobremesa ou lanche). 
Essa hierarquia é conceitual e não é implementada como tabelas separadas, apenas reflete na estrutura do conteúdo JSONB.

\begin{figure}[H]
    \centering
    \caption{Hierarquia conceitual dos itens por categoria.}
    \includegraphics[width=0.6\textwidth]{hierarquia-item.png}
    \label{fig:hierarquia-item}
    \legend{Fonte: Elaborado pela autora.}
\end{figure}

\subsubsection{Definição dos Atributos (PostgreSQL)}

A Tabela~\ref{tab:atributos-restaurante} apresenta os atributos da tabela \texttt{restaurantes}.

\begin{table}[H]
    \centering
    \renewcommand{\arraystretch}{1.1}
    \caption{Atributos da tabela \texttt{restaurantes}}
    \label{tab:atributos-restaurante}
    \small
    \begin{tabular}{|>{\raggedright\arraybackslash}p{3cm}|>{\raggedright\arraybackslash}p{3.2cm}|>{\raggedright\arraybackslash}p{8cm}|}
        \hline
        \textbf{Atributo} & \textbf{Tipo} & \textbf{Descrição} \\
        \hline
        \texttt{id} & \texttt{UUID} & Identificador único do restaurante (chave primária). \\ \hline
        \texttt{nome} & \texttt{VARCHAR(255)} & Nome do restaurante. \\ \hline
        \texttt{categoria} & \texttt{VARCHAR(100)} & Classificação da culinária do restaurante. \\ \hline
        \texttt{ativo} & \texttt{BOOLEAN} & Indica se o restaurante está em funcionamento. \\ \hline
        \texttt{criado\_em} & \texttt{TIMESTAMP} & Data de criação do registro. \\
        \hline
    \end{tabular}
    \vspace{0.4em}
\legend{Fonte: Elaborado pela autora.}
\end{table}

Já a Tabela~\ref{tab:atributos-itens} apresenta os atributos da tabela \texttt{itens}. 
Os atributos \texttt{hash} e \texttt{versao} não pertencem ao domínio do cardápio, eles são usados na lógica de inserção condicional, explicada na C2 da Seção~\ref{sec:consultas-implementadas}. 

\begin{table}[H]
    \centering
    \renewcommand{\arraystretch}{1.1}
    \caption{Atributos da tabela \texttt{itens}}
    \label{tab:atributos-itens}
    \small
    \begin{tabular}{|>{\raggedright\arraybackslash}p{3.4cm}|>{\raggedright\arraybackslash}p{3.2cm}|>{\raggedright\arraybackslash}p{7.6cm}|}
        \hline
        \textbf{Atributo} & \textbf{Tipo} & \textbf{Descrição} \\
        \hline
        \texttt{id} & \texttt{UUID} & Identificador único do item (chave primária). \\ \hline
        \texttt{restaurante\_id} & \texttt{UUID} & Restaurante a que o item pertence (chave estrangeira). \\ \hline
        \texttt{nome} & \texttt{VARCHAR(255)} & Nome do item. \\ \hline
        \texttt{categoria} & \texttt{VARCHAR(100)} & Pizza, bebida, sobremesa ou lanche. \\ \hline
        \texttt{status} & \texttt{VARCHAR(20)} & \texttt{AVAILABLE}, \texttt{UNAVAILABLE} ou \texttt{DELETED}. O último representa a exclusão lógica. \\ \hline
        \texttt{preco} & \texttt{JSONB} & Preço do item, em um objeto de estrutura fixa. \\ \hline
        \texttt{opcoes\_grupos} & \texttt{JSONB} & Grupos de opções configuráveis pelo cliente, em um array de objetos aninhados. \\ \hline
        \texttt{atributos\_variaveis} & \texttt{JSONB} & Características descritivas que variam conforme a categoria. \\ \hline
        \texttt{hash} & \texttt{INT} & Valor derivado do conteúdo do item, usado para identificar se ele mudou. \\ \hline
        \texttt{versao} & \texttt{BIGINT} & Contador de versão do item. \\ \hline
        \texttt{criado\_em} & \texttt{TIMESTAMP} & Data de criação do registro. \\ \hline
        \texttt{modificado\_em} & \texttt{TIMESTAMP} & Data da última modificação do registro. \\ \hline
    \end{tabular}
    \vspace{0.4em}
\legend{Fonte: Elaborado pela autora.}
\end{table}

\subsubsection{Atributos JSONB}

Os três atributos JSONB foram escolhidos para representar padrões estruturais distintos de dados semiestruturados, 
que se comportam de maneira diferente nas consultas e nos índices.

O atributo \texttt{preco} é um objeto simples, com estrutura igual para todos os itens (Tabela~\ref{tab:jsonb-preco}). 

\begin{table}[H]
\centering
\renewcommand{\arraystretch}{1.1}
\caption{Estrutura do atributo \texttt{preco}.}
\label{tab:jsonb-preco}
\small
\begin{tabular}{|l|l|p{6cm}|l|}
\hline
\textbf{Atributo} & \textbf{Tipo} & \textbf{Descrição} & \textbf{Obrigatório?} \\ \hline
\texttt{valor} & Number (float) & Preço de venda atual do item. & Sim \\ \hline
\texttt{moeda} & String & Código da moeda corrente. & Sim \\ \hline
\texttt{precoOriginal} & Number (float) ou null & Preço antes do desconto. É preenchido apenas quando o item possui promoção, caso contrário, é \texttt{null}. & Não \\ \hline
\end{tabular}
\vspace{0.4em}
\legend{Fonte: Elaborado pela autora.}
\end{table}

O atributo \texttt{opcoes\_grupos} é um array de objetos aninhados (Tabela~\ref{tab:jsonb-opcoes}). Cada grupo é um
objeto com os campos \texttt{nome}, \texttt{minimo}, \texttt{maximo} (limites mínimo e máximo de escolhas) e
\texttt{opcoes}. 

\begin{table}[H]
\centering
\renewcommand{\arraystretch}{1.1}
\caption{Estrutura do atributo \texttt{opcoes\_grupos}.}
\label{tab:jsonb-opcoes}
\small
\begin{tabular}{|l|l|l|p{4.2cm}|l|}
\hline
\textbf{Categoria} & \textbf{Grupo esperado} & \textbf{Limites} & \textbf{Estrutura de \texttt{opcoes}} & \textbf{Obrigatório?} \\ \hline
\textbf{pizza} & \texttt{Borda} & min: 0, max: 1 & Array de \{nome, precoAdicional\} & Não \\ \cline{2-5}
 & \texttt{Tamanho} & min: 1, max: 1 & Array de \{nome, precoAdicional\} & Sim \\ \hline
\textbf{bebida} & \texttt{Tamanho} & min: 1, max: 1 & Array de \{nome, precoAdicional\} & Sim \\ \hline
\textbf{sobremesa} & \texttt{Acompanhamento} & min: 0, max: 2 & Array de \{nome, precoAdicional\} & Não \\ \hline
\textbf{lanche} & \texttt{Adicionais} & min: 0, max: 5 & Array de \{nome, precoAdicional\} & Não \\ \hline
\end{tabular}
\vspace{0.4em}
\legend{Fonte: Elaborado pela autora.}
\end{table}

O atributo \texttt{atributos\_variaveis} é um objeto cujas chaves dependem da categoria do item 
(Tabela~\ref{tab:jsonb-atributos}). Atributos como \texttt{tags} e \texttt{promocao} são opcionais e ocorrem em mais de uma categoria. 
Neste caso, atributos não obrigatórios são omitidos quando não se aplicam, em vez de armazenados como nulos. O campo \texttt{promocao} é um objeto com \texttt{descontoPercentual} (inteiro) e \texttt{validoAte}
(texto com a data), e \texttt{informacaoNutricional} é um objeto com \texttt{calorias} (inteiro) e
\texttt{contemGluten} (booleano).

\begin{table}[H]
\centering
\renewcommand{\arraystretch}{1.1}
\caption{Estrutura do atributo \texttt{atributos\_variaveis}.}
\label{tab:jsonb-atributos}
\small
\begin{tabular}{|l|l|l|l|}
\hline
\textbf{Categoria} & \textbf{Atributo esperado} & \textbf{Tipo} & \textbf{Obrigatório?} \\ \hline
\textbf{pizza} & \texttt{promocao} & Objeto \{descontoPercentual, validoAte\} & Não \\ \cline{2-4}
 & \texttt{tags} & Array de strings & Não \\ \hline
\textbf{bebida} & \texttt{gelada} & Boolean & Sim \\ \cline{2-4}
 & \texttt{tags} & Array de strings & Não \\ \hline
\textbf{sobremesa} & \texttt{informacaoNutricional} & Objeto \{calorias, contemGluten\} & Sim \\ \cline{2-4}
 & \texttt{promocao} & Objeto \{descontoPercentual, validoAte\} & Não \\ \hline
\textbf{lanche} & \texttt{ingredientes} & Array de strings & Sim \\ \cline{2-4}
 & \texttt{tags} & Array de strings & Não \\ \cline{2-4}
 & \texttt{promocao} & Objeto \{descontoPercentual, validoAte\} & Não \\ \hline
\end{tabular}
\vspace{0.4em}
\legend{Fonte: Elaborado pela autora.}
\end{table}

\subsubsection{Criação das tabelas no PostgreSQL}

No PostgreSQL, o esquema é definido de forma explícita antes da inserção dos dados. O Código~\ref{lst:ddl} apresenta os comandos de criação das duas tabelas. A chave estrangeira \texttt{fk\_item\_restaurante} garante a integridade referencial entre item e restaurante, e a cláusula \texttt{ON DELETE CASCADE} remove os itens de um restaurante quando ele é removido.

\begin{codigo}[language=SQL, breaklines=true, caption={Criação das tabelas no PostgreSQL.}, label={lst:ddl}]
CREATE TABLE restaurantes (
  id          UUID DEFAULT gen_random_uuid(),
  nome        VARCHAR(255) NOT NULL,
  categoria   VARCHAR(100),
  ativo       BOOLEAN NOT NULL DEFAULT true,
  criado_em   TIMESTAMP NOT NULL DEFAULT now(),
  CONSTRAINT pk_restaurante PRIMARY KEY (id)
);

CREATE TABLE itens (
  id                    UUID DEFAULT gen_random_uuid(),
  restaurante_id        UUID NOT NULL,
  nome                  VARCHAR(255) NOT NULL,
  categoria             VARCHAR(100) NOT NULL,
  status                VARCHAR(20) NOT NULL DEFAULT 'AVAILABLE',
  preco                 JSONB NOT NULL,
  opcoes_grupos         JSONB,
  atributos_variaveis   JSONB NOT NULL DEFAULT '{}'::jsonb,
  hash                  INT,
  versao                BIGINT NOT NULL DEFAULT 1,
  criado_em             TIMESTAMP NOT NULL DEFAULT now(),
  modificado_em         TIMESTAMP NOT NULL DEFAULT now(),
  CONSTRAINT pk_item PRIMARY KEY (id),
  CONSTRAINT fk_item_restaurante FOREIGN KEY (restaurante_id)
    REFERENCES restaurantes(id) ON DELETE CASCADE
);
\end{codigo}

\subsubsection{Modelagem no MongoDB}

No MongoDB não há um esquema equivalente ao do PostgreSQL. As coleções \texttt{restaurantes} e \texttt{itens} não são declaradas previamente: elas passam a existir 
quando os primeiros documentos são inseridos, ou seja, a estrutura de cada documento é definida pelo próprio conteúdo inserido. Essa ausência de integridade referencial é uma 
característica dos bancos não relacionais discutida na Seção~\ref{sec:relacional-vs-nao-relacional}, ele transfere para a aplicação a responsabilidade de manter a consistência 
entre item e restaurante.

A principal diferença estrutural está nos atributos variáveis. No PostgreSQL, eles ficam dentro da coluna \texttt{atributos\_variaveis}, já no MongoDB 
eles são gravados como campos de nível superior do documento do item, ao lado de \texttt{nome} e \texttt{categoria}. Isso altera a forma de consultá-los: 
no PostgreSQL, o acesso parte da coluna \texttt{atributos\_variaveis}, e no MongoDB, ele é feito diretamente pelo nome do campo no documento. 
Já \texttt{preco} e \texttt{opcoes\_grupos} são, em ambos os bancos, um objeto e um array de objetos, respectivamente, mas no MongoDB são campos nativos do documento, 
e não valores de uma coluna. 

A Tabela~\ref{tab:comparacao-modelos} resume as diferenças entre os dois modelos.

\begin{table}[H]
    \centering
    \caption{Comparação entre os modelos do PostgreSQL e do MongoDB}
    \label{tab:comparacao-modelos}
    \small
    \begin{tabular}{|>{\raggedright\arraybackslash}p{3.6cm}|>{\raggedright\arraybackslash}p{5.4cm}|>{\raggedright\arraybackslash}p{5.4cm}|}
        \hline
        \textbf{Aspecto} & \textbf{PostgreSQL} & \textbf{MongoDB} \\
        \hline
        Esquema & Definido antes da inserção, com tipos e restrições & Definido pelo conteúdo inserido, sem validação \\ \hline
        Identificador & \texttt{id} (UUID) & \texttt{\_id} (ObjectId) \\ \hline
        Referência ao restaurante & \texttt{restaurante\_id}, com chave estrangeira & \texttt{restauranteId}, sem integridade referencial \\ \hline
        Preço e opções & Colunas JSONB & Objeto e array embutidos no documento \\ \hline
        Atributos variáveis & Coluna \texttt{atributos\_variaveis} & Campos de nível superior do documento \\ \hline
    \end{tabular}
    \vspace{0.4em}
\legend{Fonte: Elaborado pela autora.}
\end{table}

\subsubsection{Limpeza do banco}

Antes de cada carga de dados, os dois bancos são limpos para que execuções anteriores não interfiram nos resultados. No PostgreSQL, as tabelas são removidas com \texttt{DROP TABLE IF EXISTS ... CASCADE} e recriadas com os comandos do Código~\ref{lst:ddl}. No MongoDB, as duas coleções são removidas com \texttt{drop()} e voltam a existir na inserção seguinte. 

\subsection{Geração dos Dados}

Como não foi possível utilizar dados reais do catálogo que inspira este trabalho, os dados do experimento são gerados de forma sintética, por meio de um \textit{script} em Python. A geração usa duas bibliotecas: a biblioteca padrão \texttt{random}, que realiza os sorteios, e a biblioteca \texttt{Faker}, configurada para o português do Brasil, que fornece sobrenomes e nomes de cidades para compor os nomes dos restaurantes. Os demais textos, como os nomes dos itens, dos ingredientes e das opções, vêm de listas definidas pela autora. As probabilidades descritas nesta seção também foram definidas de forma arbitrária, com o objetivo de tornar o catálogo heterogêneo.

Os mesmos dados gerados são inseridos nos dois bancos, de modo que a comparação ocorre sobre conteúdos
idênticos. Apenas o formato do identificador muda: UUID no PostgreSQL e \texttt{ObjectId} no MongoDB.
As sementes dos geradores aleatórios (\texttt{random.seed(42)} e \texttt{Faker.seed(42)}) são fixas, de
modo que o conteúdo gerado se repete a cada execução. Os identificadores, por outro lado, não dependem
dessas sementes e mudam a cada execução.


\subsubsection{Restaurantes}

A Tabela~\ref{tab:geracao-restaurantes} resume como cada atributo do restaurante é gerado.

\begin{table}[H]
    \centering
    \renewcommand{\arraystretch}{1.1}
    \caption{Geração dos atributos dos restaurantes}
    \label{tab:geracao-restaurantes}
    \small
    \begin{tabular}{|>{\raggedright\arraybackslash}p{3cm}|>{\raggedright\arraybackslash}p{6cm}|>{\raggedright\arraybackslash}p{5.2cm}|}
        \hline
        \textbf{Atributo} & \textbf{Geração} & \textbf{Distribuição} \\ \hline
        Categoria & Atribuída de forma cíclica & Italiana, brasileira, árabe e mexicana, em partes iguais \\ \hline
        Nome & Prefixo (três por categoria, como \textit{Cantina}, \textit{Trattoria} e \textit{Pizzaria}) mais sobrenome ou nome de cidade do \texttt{Faker} & Único em todo o conjunto \\ \hline
        Porte & Sorteio entre três portes, que define o peso do restaurante na distribuição dos itens & Pequeno (60\%), médio (30\%) e grande (10\%) \\ \hline
    \end{tabular}
    \vspace{0.4em}
\legend{Fonte: Elaborado pela autora.}
\end{table}

Os pesos de porte fazem com que os restaurantes grandes concentrem,
em média, cerca de 49\% dos itens, os médios cerca de 37\% e os pequenos cerca de 15\%. Essa desigualdade
é proposital, pois permite medir como a seletividade afeta uma consulta.

\subsubsection{Itens}

A Tabela~\ref{tab:geracao-itens} resume como cada atributo de Item é gerado.

\begin{table}[H]
    \centering
    \renewcommand{\arraystretch}{1.1}
    \caption{Geração dos atributos dos itens}
    \label{tab:geracao-itens}
    \small
    \begin{tabular}{|>{\raggedright\arraybackslash}p{3cm}|>{\raggedright\arraybackslash}p{6cm}|>{\raggedright\arraybackslash}p{5.2cm}|}
        \hline
        \textbf{Atributo} & \textbf{Geração} & \textbf{Distribuição} \\ \hline
        Categoria & Sorteio com probabilidades distintas & Pizza: 35\% \newline Bebida: 20\% \newline Sobremesa: 10\% \newline Lanche: 35\% \\ \hline
        Nome & Nome-base (seis a oito por categoria, como \textit{Margherita}, \textit{Brownie} e \textit{X-Bacon}) mais um complemento opcional (seis complementos possíveis ou nenhum, como \textit{Tradicional} e \textit{da Casa}) & Não precisa ser único \\ \hline
        Status & Sorteio entre três valores & \texttt{AVAILABLE} (85\%) \newline \texttt{UNAVAILABLE} (10\%) \newline e \texttt{DELETED} (5\%) \\ \hline
        Atributos variáveis & Definidos por categoria, com sorteios independentes & Tabela~\ref{tab:probabilidades} \\ \hline
        Opções de grupos & Definidas por categoria, com preço adicional fixo por opção & Tabela~\ref{tab:opcoes-geradas} \\ \hline
        Preço & Valor uniforme na faixa da categoria, com duas casas decimais & Pizza: R\$~35 a R\$~70 \newline Bebida: R\$~5 a R\$~15 \newline Sobremesa: R\$~12 a R\$~25 \newline Lanche: R\$~18 a R\$~35 \\ \hline
    \end{tabular}
    \vspace{0.4em}
\legend{Fonte: Elaborado pela autora.}
\end{table}

A distribuição das categorias foi definida para que o catálogo tivesse categorias de tamanhos diferentes, como ocorre em catálogos reais. 

Os atributos de \texttt{atributos\_variaveis} são definidos por categoria e cada atributo 
opcional é decidido por um sorteio independente. A lógica é: o gerador compara um número 
aleatório entre 0 e 1 com um limiar, o atributo só é incluído quando o número do sorteio 
o supera. Atributos que não se aplicam são omitidos do objeto. A Tabela~\ref{tab:probabilidades} 
resume os atributos gerados e suas probabilidades.

\begin{table}[H]
    \centering
    \renewcommand{\arraystretch}{1.1}
    \caption{Atributos variáveis gerados por categoria}
    \label{tab:probabilidades}
    \small
    \begin{tabular}{|>{\raggedright\arraybackslash}p{3.3cm}|>{\raggedright\arraybackslash}p{2.7cm}|>{\raggedright\arraybackslash}p{3.1cm}|>{\raggedright\arraybackslash}p{5.2cm}|}
        \hline
        \textbf{Atributo} & \textbf{Categorias} & \textbf{Presença} & \textbf{Conteúdo gerado} \\ \hline
        \texttt{tags} & Pizza \newline Bebida \newline Lanche & Opcional & Sorteadas entre sete tags possíveis, cada uma com 30\% de chance, então um item pode ter várias tags. Se nenhuma for sorteada, o atributo é omitido. \\ \hline
        \texttt{promocao} & Pizza \newline Sobremesa \newline Lanche & Opcional: \newline 40\% \newline 30\% \newline 35\% & Objeto com \texttt{descontoPercentual} (inteiro entre 5 e 30 em 98\% das promoções, e 50, 60, 70 ou 80 nos 2\% restantes) e \texttt{validoAte} (data fixa). \\ \hline
        \texttt{gelada} & Bebida & Sempre & Booleano, verdadeiro em 80\% dos casos. \\ \hline
        \texttt{informacao\allowbreak Nutricional} & Sobremesa & Sempre & Objeto com \texttt{calorias} (inteiro entre 200 e 500) e \texttt{contemGluten} (verdadeiro em 50\% dos casos). \\ \hline
        \texttt{ingredientes} & Lanche & Sempre & Lista com dois ou mais ingredientes, sorteados entre dez possíveis, cada um com 50\% de chance, então um lanche pode ter vários ingredientes. Se saírem menos de dois, usa \textit{queijo} e \textit{bacon}. \\ \hline
    \end{tabular}
    \vspace{0.4em}
\legend{Fonte: Elaborado pela autora.}
\end{table}

Os grupos de \texttt{opcoes\_grupos} são definidos por categoria, com os limites apresentados na Tabela~\ref{tab:jsonb-opcoes}. A Tabela~\ref{tab:opcoes-geradas} apresenta as opções de cada grupo e o preço adicional fixo de cada uma.

\begin{table}[H]
    \centering
    \renewcommand{\arraystretch}{1.1}
    \caption{Opções geradas em cada grupo}
    \label{tab:opcoes-geradas}
    \small
    \begin{tabular}{|l|l|l|l|>{\raggedright\arraybackslash}p{3.6cm}|}
        \hline
        \textbf{Categoria} & \textbf{Grupo} & \textbf{Opção} & \textbf{Preço adicional} & \textbf{Presença} \\ \hline
        \textbf{pizza} & \texttt{Borda} & Sem borda & R\$~0,00 & Sempre \\ \cline{3-5}
         & & Borda mussarela & R\$~8,00 & Sempre \\ \cline{3-5}
         & & Borda catupiry & R\$~10,00 & Sempre \\ \cline{2-5}
         & \texttt{Tamanho} & Média & R\$~0,00 & Sempre \\ \cline{3-5}
         & & Grande & R\$~12,00 & Sempre \\ \cline{3-5}
         & & Família & R\$~20,00 & Sempre \\ \hline
        \textbf{bebida} & \texttt{Tamanho} & 350ml & R\$~0,00 & Sempre \\ \cline{3-5}
         & & 500ml & R\$~2,00 & Sempre \\ \cline{3-5}
         & & 600ml & R\$~3,50 & Sempre \\ \hline
        \textbf{sobremesa} & \texttt{Acompanhamento} & Calda & R\$~2,00 & Sempre \\ \cline{3-5}
         & & Sem acompanhamento & R\$~0,00 & Sempre \\ \hline
        \textbf{lanche} & \texttt{Adicionais} & Sem cebola & R\$~0,00 & Sempre \\ \cline{3-5}
         & & Bacon extra & R\$~3,00 & Restaurantes que o oferecem (cerca de 60\%) \\ \cline{3-5}
         & & Cheddar extra & R\$~3,50 & Restaurantes que o oferecem (cerca de 25\%) \\ \cline{3-5}
         & & Picles extra & R\$~2,50 & Restaurantes que o oferecem (cerca de 1\%) \\ \hline
    \end{tabular}
    \vspace{0.4em}
\legend{Fonte: Elaborado pela autora.}
\end{table}

A oferta de \textit{Bacon extra}, \textit{Cheddar extra} e \textit{Picles extra} é decidida no nível do restaurante e não do item, porque não faria sentido um mesmo restaurante oferecê-los em alguns lanches e em outros não. Dessa forma, a diferença de oferta entre os restaurantes produz seletividades distintas nas consultas que filtram por essas opções. 

O preço é um objeto, cujo campo \texttt{precoOriginal} só é preenchido quando o item possui promoção e é calculado a partir do percentual de desconto. Por esse motivo, os atributos variáveis são gerados antes do preço.

Cada item recebe um valor de \texttt{hash}, que funciona como uma impressão digital do seu conteúdo. Ele é calculado a partir do nome, da categoria, do status, do preço, das opções e dos atributos variáveis, serializados em JSON com as chaves ordenadas, o que garante que o mesmo conteúdo produza sempre o mesmo valor, o Código~\ref{lst:hash}, no Apêndice~\ref{chapter:apendice-hash}, apresenta esse cálculo. O objetivo é permitir, na inserção condicional, decidir se um item recebido difere do armazenado comparando apenas dois inteiros, em vez de comparar todos os campos.

\subsection{Inserção dos dados}

Os restaurantes são inseridos antes dos itens, porque os itens referenciam o restaurante. 
Os dois bancos recebem os itens em lotes de 10000, e cada lote é confirmado da seguinte forma: 
no PostgreSQL, com um \texttt{commit} ao final de cada lote, e no MongoDB com a opção 
\texttt{journal=True}. 
O PostgreSQL, gera o identificador de cada 
item (\texttt{DEFAULT gen\_random\_uuid()}) e preenche as datas de criação e modificação 
(\texttt{DEFAULT now()}). O Código~\ref{lst:insert-pg}, no Apêndice~\ref{chapter:apendice-codigos}, apresenta a inserção no PostgreSQL.

No MongoDB, os itens são carregados com uma chamada a \texttt{insert\_many} por lote, e o identificador do
item é gerado pelo \textit{driver}. Como o banco não preenche as datas, a aplicação as define com o instante
do início da carga, o mesmo valor para todos os documentos (Código~\ref{lst:insert-mongo}, Apêndice~\ref{chapter:apendice-codigos}).


Durante cada inserção são medidos o tempo total, o uso de CPU e de memória do contêiner e o 
espaço ocupado pelos dados. Ao final da carga, o PostgreSQL executa o comando
\texttt{ANALYZE} nas duas tabelas, para que todos os cenários partam das mesmas estatísticas, em vez de
depender do momento em que o \textit{autovacuum} rodaria sozinho.

\subsection{Consultas Implementadas}
\label{sec:consultas-implementadas}

As consultas do experimento foram definidas a partir de dois critérios. O primeiro é a 
recorrência: foram escolhidas operações que ocorrem com frequência no cenário do iFood, 
como carregar o cardápio de um restaurante, filtrar itens por categoria, processar eventos de 
atualização do catálogo e remover itens do cardápio. O segundo é a investigação de possíveis 
gargalos: como a busca em arrays de objetos aninhados dentro de um campo JSONB e a 
busca por um valor específico dentro de um atributo opcional, 
como o valor da promoção, presente em apenas parte dos itens. Cada consulta 
representa um cenário distinto de uso do catálogo.

As oito consultas foram divididas em dois grupos, conforme o tipo de dado sobre o qual atuam: 
consultas sobre dados estruturados (C1 a C4) que operam sobre campos escalares comuns e
consultas sobre dados semiestruturados (C5 a C8) que operam sobre o conteúdo dos atributos 
variáveis. A Tabela~\ref{tab:consultas} resume as consultas, e as seções seguintes apresentam 
cada uma com a implementação nos dois bancos. 

\begin{table}[H]
    \centering
    \renewcommand{\arraystretch}{1.1}
    \caption{Resumo das consultas implementadas}
    \label{tab:consultas}
    \small
    \begin{tabular}{|l|>{\raggedright\arraybackslash}p{2.5cm}|>{\raggedright\arraybackslash}p{1.8cm}|>{\raggedright\arraybackslash}p{3.9cm}|>{\raggedright\arraybackslash}p{3.9cm}|}
        \hline
        \textbf{Id.} & \textbf{Grupo} & \textbf{Operação} & \textbf{Padrão de acesso} & \textbf{Cenário representado} \\ \hline
        C1 & Estruturado & Leitura & Igualdade na categoria e exclusão de um status & Listar os itens não excluídos de uma categoria específica. \\ \hline
        C2 & Estruturado & Escrita & Upsert condicional por identificador & Processar um evento de atualização do catálogo. \\ \hline
        C3 & Estruturado & Escrita & \textit{Soft delete} por identificador & Remover um item sem excluí-lo fisicamente. \\ \hline
        C4 & Estruturado & Junção & Junção e filtro de status & Carregar o cardápio completo de um restaurante. \\ \hline
        C5 & Semiestruturado & Leitura & Intervalo em campo numérico do JSONB & Filtrar itens por faixa de preço. \\ \hline
        C6 & Semiestruturado & Leitura & Igualdade em valor de atributo aninhado & Buscar itens em promoção. \\ \hline
        C7 & Semiestruturado & Leitura & Contenção em array aninhado & Buscar itens que oferecem uma opção. \\ \hline
        C8 & Semiestruturado & Escrita & Atualização parcial do JSONB por identificador & Adicionar uma tag a um item. \\ \hline
    \end{tabular}
    \vspace{0.4em}
\legend{Fonte: Elaborado pela autora.}
\end{table}

\subsubsection{C1: itens por categoria e status.} Retorna os itens de uma categoria, por exemplo \textit{pizza}, que não tenham sido excluídos logicamente. Essa consulta representa o cliente que busca um tipo de produto e recebe os itens dessa categoria. O padrão de acesso é a igualdade em \texttt{categoria} combinada com a exclusão de um valor de \texttt{status}, ambos campos escalares. O Código~\ref{lst:c1-pg} apresenta a consulta no PostgreSQL, e o Código~\ref{lst:c1-mongo}, no MongoDB.


\begin{codigo}[language=SQL, breaklines=true, caption={C1 no PostgreSQL.}, label={lst:c1-pg}]
SELECT nome, categoria, status
FROM itens
WHERE categoria = :categoria AND status <> 'DELETED';
\end{codigo}

\begin{codigo}[language=Python, breaklines=true, caption={C1 no MongoDB.}, label={lst:c1-mongo}]
db.itens.find(
    {"categoria": categoria, "status": {"$ne": "DELETED"}},
    {"_id": 0, "nome": 1, "categoria": 1, "status": 1}
)
\end{codigo}

\subsubsection{C2: upsert condicional.} Representa a operação de escrita mais frequente de 
catálogo do iFood. Insere um item novo ou atualiza um já existente, aplicando a mudança 
apenas se o conteúdo realmente mudou (\texttt{hash} diferente) e se o evento não é mais 
antigo que o já processado (\texttt{versao} maior ou igual). No PostgreSQL, a decisão é tomada 
pelo próprio banco, com
um \texttt{UPDATE} condicionado a essas duas regras, seguido de um \texttt{INSERT} que só 
ocorre se o item ainda não existe (Código~\ref{lst:c2-pg}). 
Não foi usado o comando \texttt{INSERT ... ON CONFLICT}, que seria a forma mais direta e eficiente, 
pois ele sempre recorre ao índice da chave 
primária para detectar o conflito, o que impediria medir a operação sem índice, cenário 
necessário para a análise apresentada adiante.

\begin{codigo}[language=SQL, breaklines=true, caption={C2 no PostgreSQL.}, label={lst:c2-pg}]
WITH atualizado AS (
  UPDATE itens
  SET nome = :nome, categoria = :categoria, status = :status,
      preco = CAST(:preco AS JSONB),
      opcoes_grupos = CAST(:opcoes_grupos AS JSONB),
      atributos_variaveis = CAST(:atributos AS JSONB),
      hash = :hash, versao = :versao, modificado_em = now()
  WHERE id = CAST(:id AS UUID) AND hash <> :hash AND :versao >= versao
  RETURNING id
)
INSERT INTO itens (id, restaurante_id, nome, categoria, status, preco, opcoes_grupos, atributos_variaveis, hash, versao)
SELECT CAST(:id AS UUID), CAST(:restaurante_id AS UUID), :nome, :categoria, :status, CAST(:preco AS JSONB), CAST(:opcoes_grupos AS JSONB), CAST(:atributos AS JSONB), :hash, :versao
WHERE NOT EXISTS (SELECT 1 FROM itens WHERE id = CAST(:id AS UUID));
\end{codigo}

No MongoDB, a mesma lógica é executada pela aplicação, que lê primeiro o documento existente 
e só então decide entre inserir, substituir ou não fazer nada (Código~\ref{lst:c2-mongo}).


\begin{codigo}[language=Python, breaklines=true, caption={C2 no MongoDB.}, label={lst:c2-mongo}]
existente = db.itens.find_one({"_id": evento["_id"]})
if existente is None:
    db.itens.insert_one(evento)
elif existente["hash"] != evento["hash"] and evento["versao"] >= existente["versao"]:
    evento["criado_em"] = existente["criado_em"] 
    db.itens.replace_one({"_id": evento["_id"]}, evento)
# caso contrario, nao faz nada
\end{codigo}

\subsubsection{C3: exclusão lógica.} Representa a remoção de um item do cardápio, o sistema real não executa exclusões físicas, apenas lógicas. O padrão de acesso é a atualização do campo \texttt{status}, localizando o item pelo identificador. Além do \texttt{status}, a operação zera o \texttt{hash} e atualiza \texttt{modificado\_em}. Os Códigos~\ref{lst:c3-pg} e~\ref{lst:c3-mongo} apresentam a operação nos dois bancos.

\begin{codigo}[language=SQL, breaklines=true, caption={C3 no PostgreSQL.}, label={lst:c3-pg}]
UPDATE itens
SET status = 'DELETED', hash = 0, modificado_em = now()
WHERE id = ANY(CAST(:ids AS uuid[])) AND status <> 'DELETED';
\end{codigo}

\begin{codigo}[language=Python, breaklines=true, caption={C3 no MongoDB.}, label={lst:c3-mongo}]
db.itens.update_many(
    {"_id": {"$in": ids}, "status": {"$ne": "DELETED"}},
    {"$set": {"status": "DELETED", "hash": 0, "modificado_em": datetime.now()}}
)
\end{codigo}

\subsubsection{C4: itens de um restaurante.} Retorna os itens de um restaurante específico, 
excluindo os removidos. O padrão de acesso é uma junção entre itens e restaurantes, 
com filtro de status, e a busca é feita pelo nome do restaurante. No PostgreSQL, 
a junção é resolvida pelo banco em uma única instrução (Código~\ref{lst:c4-pg}). 

\begin{codigo}[language=SQL, breaklines=true, caption={C4 no PostgreSQL.}, label={lst:c4-pg}]
SELECT i.nome, i.categoria, i.status, r.nome AS restaurante_nome
FROM itens i
JOIN restaurantes r ON r.id = i.restaurante_id
WHERE r.nome = :nome AND i.status <> 'DELETED';
\end{codigo}

No MongoDB, a aplicação busca primeiro o restaurante pelo nome, para obter seu identificador, e em seguida consulta os itens que o referenciam (Código~\ref{lst:c4-mongo}). Embora o MongoDB ofereça o operador \texttt{\$lookup}, optou-se por duas consultas simples, que é a forma usual de resolver uma referência entre coleções. 

\begin{codigo}[language=Python, breaklines=true, caption={C4 no MongoDB.}, label={lst:c4-mongo}]
restaurante = db.restaurantes.find_one({"nome": nome})
db.itens.find(
    {"restauranteId": restaurante["_id"], "status": {"$ne": "DELETED"}},
    {"_id": 0, "nome": 1, "categoria": 1, "status": 1}
)
\end{codigo}


\subsubsection{C5: itens por faixa de preço.} Retorna os itens cujo preço está entre um valor mínimo e um máximo. Como o preço é armazenado dentro do atributo JSONB, o filtro precisa acessar o campo \texttt{valor} dentro do objeto. No PostgreSQL, o valor é extraído como texto com o operador \texttt{->>} e convertido para numérico (Código~\ref{lst:c5-pg}).

\begin{codigo}[language=SQL, breaklines=true, caption={C5 no PostgreSQL.}, label={lst:c5-pg}]
SELECT nome, categoria, preco
FROM itens
WHERE (preco->>'valor')::numeric BETWEEN :minimo AND :maximo
  AND status <> 'DELETED';
\end{codigo}

No MongoDB, o campo interno é acessado diretamente por notação de ponto (Código~\ref{lst:c5-mongo}).

\begin{codigo}[language=Python, breaklines=true, caption={C5 no MongoDB.}, label={lst:c5-mongo}]
db.itens.find(
    {"preco.valor": {"$gte": minimo, "$lte": maximo},
     "status": {"$ne": "DELETED"}},
    {"_id": 0, "nome": 1, "categoria": 1, "preco": 1}
)
\end{codigo}

\subsubsection{C6: valor de atributo variável aninhado.} Representa a busca por itens com 
uma característica rara dentro de um atributo opcional, como as promoções de desconto alto. 
O padrão de acesso é a igualdade em um valor aninhado:  
\texttt{promocao.descontoPercentual} igual a 80\%. 
No PostgreSQL, usa-se o operador de contenção \texttt{@>} sobre o JSONB 
(Código~\ref{lst:c6-pg}), e no MongoDB, a notação de ponto sobre o campo aninhado 
(Código~\ref{lst:c6-mongo}). 

\begin{codigo}[language=SQL, breaklines=true, caption={C6 no PostgreSQL.}, label={lst:c6-pg}]
SELECT nome, categoria, atributos_variaveis->'promocao' AS promocao
FROM itens
WHERE atributos_variaveis @> CAST(:filtro AS JSONB)
  AND status <> 'DELETED';
-- filtro: {"promocao": {"descontoPercentual": 80}}
\end{codigo}

\begin{codigo}[language=Python, breaklines=true, caption={C6 no MongoDB.}, label={lst:c6-mongo}]
db.itens.find(
    {"promocao.descontoPercentual": 80, "status": {"$ne": "DELETED"}},
    {"_id": 0, "nome": 1, "categoria": 1, "promocao": 1}
)
\end{codigo}


\subsubsection{C7: filtro em array de objetos aninhados.} Representa a busca por itens que oferecem uma personalização específica, como \textit{Bacon extra}, e envolve a estrutura mais complexa do modelo: um array de grupos, cada um com seu próprio array de opções, o que resulta em dois níveis de aninhamento. No PostgreSQL, usa-se o operador de contenção \texttt{@>} com um filtro JSONB que descreve a estrutura procurada (Código~\ref{lst:c7-pg}). 

\begin{codigo}[language=SQL, breaklines=true, caption={C7 no PostgreSQL.}, label={lst:c7-pg}]
SELECT nome, categoria, opcoes_grupos
FROM itens
WHERE opcoes_grupos @> CAST(:filtro AS JSONB) AND status <> 'DELETED';
-- filtro: [{"opcoes": [{"nome": opcao}]}]
\end{codigo}

No MongoDB, a notação de ponto atravessa os dois arrays e compara o nome da opção diretamente (Código~\ref{lst:c7-mongo}).

\begin{codigo}[language=Python, breaklines=true, caption={C7 no MongoDB.}, label={lst:c7-mongo}]
db.itens.find(
    {"opcoes_grupos.opcoes.nome": opcao,
     "status": {"$ne": "DELETED"}},
    {"_id": 0, "nome": 1, "categoria": 1, "opcoes_grupos": 1}
)
\end{codigo}

\subsubsection{C8: atualização parcial de JSONB.} Adiciona a tag \textit{novidade} a um item específico. O padrão de acesso é a atualização de um campo interno do JSONB. No PostgreSQL, a tag é acrescentada à lista existente com \texttt{jsonb\_set} e concatenação, usando \texttt{COALESCE} para o caso de o item ainda não ter tags. Como o operador \texttt{||} apenas concatena e duplicaria a tag em um item que já a possui, a cláusula \texttt{WHERE} seleciona somente os itens que ainda não contêm o valor (Código~\ref{lst:c8-pg}).

\begin{codigo}[language=SQL, breaklines=true, caption={C8 no PostgreSQL.}, label={lst:c8-pg}]
UPDATE itens
SET atributos_variaveis = jsonb_set(
      atributos_variaveis, '{tags}',
      COALESCE(atributos_variaveis->'tags', '[]'::jsonb) || '"novidade"'::jsonb
    ),
    modificado_em = now()
WHERE id = ANY(CAST(:ids AS uuid[]))
  AND NOT COALESCE(atributos_variaveis->'tags', '[]'::jsonb) @> '"novidade"'::jsonb;
\end{codigo}

No MongoDB, o operador \texttt{\$addToSet} já evita duplicatas. Ainda assim, o filtro \texttt{"tags": \{"\$ne": "novidade"\}} foi incluído para que ambas as consultas atualizem exatamente o mesmo conjunto de itens, isto é, apenas os que ainda não possuem a tag (Código~\ref{lst:c8-mongo}). Desse modo, em ambos os bancos, a data de modificação só é alterada quando há mudança real.

\begin{codigo}[language=Python, breaklines=true, caption={C8 no MongoDB.}, label={lst:c8-mongo}]
db.itens.update_many(
    {"_id": {"$in": ids}, "tags": {"$ne": "novidade"}},
    {"$addToSet": {"tags": "novidade"},
     "$set": {"modificado_em": datetime.now()}}
)
\end{codigo}

\subsection{Estratégias de Indexação}
\label{sec:indexacao}

Esta seção descreve as estratégias de indexação avaliadas para cada consulta, 
cada uma foi executada primeiro
sem índices adicionais e depois com os índices candidatos. Os
índices foram escolhidos conforme o padrão de acesso de cada consulta
e os incompatíveis com esse padrão não foram testados. O procedimento 
é o mesmo para as duas tecnologias:

\begin{enumerate}
  \item Antes de cada configuração, todos os índices que foram anteriormente criados pela autora são
  removidos. No PostgreSQL, é usada a função \texttt{apagar\_indices\_pg()},
  que automatiza o comando \texttt{DROP INDEX}. No MongoDB, é usado
  \texttt{drop\_indexes()}. Assim, nenhum índice de uma configuração anterior 
  interfere na seguinte.
  \item Logo em seguida o índice é criado. No PostgreSQL, a criação é seguida do comando \texttt{ANALYZE}, que atualiza as estatísticas da tabela, de modo que o plano escolhido reflita o índice testado e não a desatualização das estatísticas.
  \item O plano de execução é consultado com
  \texttt{EXPLAIN} no PostgreSQL e com \texttt{explain}
  (modo \texttt{executionStats}) no MongoDB. Esse passo confirma se o índice
  criado foi escolhido pelo otimizador.
  \item Quando o planejador do PostgreSQL ignora o índice criado e opta por uma varredura sequencial, a consulta é medida também com o uso do índice forçado (\texttt{enable\_seqscan = off}), para separar o efeito do índice da decisão do planejador.
  \item Fazemos duas medições sobre os índices, a primeira é em relação ao tempo de execução da consulta e a segunda ao espaço ocupado pelo índice. O procedimento de cada uma está descrito na
  Seção~\ref{sec:metodologia-medicao}.

\end{enumerate}


A Tabela~\ref{tab:matriz-indices} resume as combinações de consulta e índice testadas.

\begin{table}[H]
    \centering
    \renewcommand{\arraystretch}{1.1}
    \caption{Combinações de consulta e índice avaliadas.}
    \label{tab:matriz-indices}
    \footnotesize
    \setlength{\tabcolsep}{4pt}
    \begin{tabular}{|>{\raggedright\arraybackslash}p{3.7cm}|>{\centering\arraybackslash}p{1.3cm}|>{\centering\arraybackslash}p{1.3cm}|>{\centering\arraybackslash}p{1.3cm}|>{\centering\arraybackslash}p{1.3cm}|>{\centering\arraybackslash}p{1.3cm}|>{\centering\arraybackslash}p{1.3cm}|>{\centering\arraybackslash}p{1.3cm}|}
        \hline
         & \multicolumn{3}{c|}{\textbf{PostgreSQL}} & \multicolumn{4}{c|}{\textbf{MongoDB}} \\ \cline{2-8}
        \textbf{Consulta} & \textbf{Sem índice} & \textbf{Simples/ Composto} & \textbf{GIN} & \textbf{Sem índice} & \textbf{Simples/ Composto} & \textbf{Multikey} & \textbf{Wildcard} \\ \hline
        C1 --- Categoria e status & \checkmark & \checkmark & -- & \checkmark & \checkmark & -- & -- \\ \hline
        C2 --- Upsert condicional & \checkmark$^{b}$ & \checkmark$^{a}$ & -- & \checkmark$^{b}$ & \checkmark$^{a}$ & -- & -- \\ \hline
        C3 --- Exclusão lógica & \checkmark$^{b}$ & \checkmark$^{a}$ & -- & \checkmark$^{b}$ & \checkmark$^{a}$ & -- & -- \\ \hline
        C4 --- Itens de um restaurante & \checkmark & \checkmark & -- & \checkmark & \checkmark & -- & -- \\ \hline
        C5 --- Faixa de preço & \checkmark & \checkmark & -- & \checkmark & \checkmark & -- & -- \\ \hline
        C6 --- Valor de atributo aninhado & \checkmark & -- & \checkmark & \checkmark & -- & -- & \checkmark \\ \hline
        C7 --- Array de objetos aninhados & \checkmark & -- & \checkmark & \checkmark & -- & \checkmark & -- \\ \hline
        C8 --- Atualização parcial & \checkmark$^{b}$ & \checkmark$^{a}$ & -- & \checkmark$^{b}$ & \checkmark$^{a}$ & -- & -- \\ \hline
    \end{tabular}
    \vspace{0.4em}
    \legend{$^{a}$ Utiliza apenas o índice da chave primária (criado automaticamente); $^{b}$ Varredura forçada da tabela; \checkmark: configuração testada; --: configuração não aplicável ao padrão de acesso. Fonte: Elaborado pela autora.}
\end{table}

A escolha de cada tipo de índice, e o motivo de não usar os outros, foi a seguinte:

\begin{itemize}
  \item \textbf{B-tree (simples ou composto)}: usado em C1, C4 e C5, que
  comparam valores escalares por igualdade ou intervalo, e em C2, C3 e C8, que utilizam
  apenas o índice da chave primária, criado automaticamente, para localizar o item pelo identificador. Em C6 e C7, as consultas usam o operador de contenção \texttt{@>} do JSONB, que o índice \textit{B-tree} não suporta.

  \item \textbf{GIN}: usado em C6 e C7, que procuram dentro do JSONB. Não é usado
  de C1 a C4, pois os campos são escalares e não têm elementos internos a
  indexar. Também não é usado em C5, pois o GIN não atende a comparações de intervalo.

  \item \textbf{\textit{Multikey}}: usado em C7, sobre o caminho
  \texttt{opcoes\_grupos.opcoes.nome}, que atravessa arrays. Não é usado no restante pois os campos consultados não são arrays.

  \item \textbf{\textit{Wildcard}}: usado em C6, em que o atributo consultado
  existe só em parte dos itens e um único índice indexa automaticamente todos os caminhos presentes em cada documento, dispensando a criação de um índice por atributo. Não é usado
  no restante, pois o campo consultado é sempre conhecido. 
\end{itemize}


\subsection{Metodologia de Medição}
\label{sec:metodologia-medicao}

As medições foram implementadas em três funções independentes, cada uma
dedicada a um certo tipo de desempenho: tempo de execução, consumo de
recursos computacionais e espaço ocupado em disco. As mesmas funções foram
aplicadas aos dois bancos, de modo que as diferenças observadas decorram do
banco e não do método de medição. 

\subsubsection{Medição de tempo}
\label{subsec:medicao-tempo}

Inicialmente, o tempo seria obtido a partir do próprio plano de execução, por
meio do \texttt{EXPLAIN ANALYZE} no PostgreSQL e do \texttt{explain} no MongoDB.
Porém, essa abordagem se mostrou inadequada para as operações de escrita, já que no
PostgreSQL o \texttt{EXPLAIN ANALYZE} executa o comando, enquanto no MongoDB o
\texttt{explain} apenas simula a operação, sem
aplicá-la. Logo, os tempos de escrita dos
dois bancos não seriam comparados de maneira justa.

Optou-se então por cronometrar a operação do lado do cliente, com a função
\texttt{time.perf\_counter()} da biblioteca padrão do Python. 
A mesma função e os mesmos pontos de início e fim
são usados para consultas, inserções, atualizações e exclusões nos dois
bancos. 

A função \texttt{medir\_tempo}, apresentada no
Código~\ref{lst:medir-tempo} (Apêndice~\ref{chapter:apendice-medicao}), recebe a ação a ser cronometrada e a executa
repetidamente. Seus parâmetros e etapas são os seguintes:

\begin{itemize}
  \item \texttt{acao}: função que executa a operação a ser
  medida (uma consulta, um \textit{insert}, um \textit{update}).
  \item \texttt{warm\_up}: quando ativo, a ação é executada uma vez antes da
  medição e o resultado é descartado. Isso evita que o custo da primeira
  execução distorça as medidas. Nos experimentos deste trabalho esse parâmetro permanece desativado,
  pois o objetivo é medir justamente a execução com o cache vazio. Mas ele pode ser ativado
  em experimentos futuros para avaliar o impacto do aquecimento na performance.
  \item \texttt{antes}: função opcional executada antes de cada repetição e
  fora da cronometragem, usada quando é preciso restaurar o estado do banco
  entre as execuções.
  \item \texttt{repeats}: número de repetições cronometradas.
\end{itemize}

Para cada ação, a função devolve o tempo mínimo, o máximo, a média e a mediana
em milissegundos, além dos tempos brutos de cada repetição. 

\subsubsection{Medição de recursos}
\label{subsec:medicao-recursos}

O consumo de CPU e de memória dos contêineres foi medido durante a inserção em massa dos dados. A leitura é feita diretamente nos contadores do cgroup do Linux (cgroup v2) do contêiner, que fornecem os valores instantaneamente. O comando \texttt{docker stats} foi mantido como alternativa, usada apenas quando o cgroup não é encontrado. 
O Código~\ref{lst:medir-recursos}, no Apêndice~\ref{chapter:apendice-recursos}, apresenta as funções que trabalham em conjunto para capturar e monitorar esses recursos:

\begin{itemize}
  \item \texttt{ler\_cgroup}: lê o tempo de CPU acumulado pelo contêiner e a 
  memória em uso. Da memória é descontado o cache de arquivos inativo, 
  como faz o \texttt{docker stats}, de modo que o valor reflita a memória efetivamente utilizada.
  \item \texttt{monitorar\_recursos}: a cada intervalo (0,25 s), calcula o percentual de 
  CPU, expresso em 
  porcentagem de um núcleo, e registra a memória, que já vem em MiB. Cada amostra, composta pelo par CPU e memória, só é 
  armazenada se o envio de dados ao banco estiver em andamento, 
  o que exclui o tempo gasto na conversão dos dados no Python.
  \item \texttt{medir\_insercao\_com\_recursos}: dispara o monitor em uma \textit{thread} 
  separada, executa a inserção na \textit{thread} principal, cronometra o tempo com o 
  \texttt{perf\_counter} e, ao final, sinaliza a parada do monitor. Devolve dois tempos, 
  o de envio ao banco (\texttt{tempo\_envio\_ms}) e o total, que inclui a conversão dos dados 
  lote a lote (\texttt{tempo\_total\_ms}), além da média e do máximo de CPU e de memória 
  entre as amostras e das leituras brutas.
\end{itemize}

Na prática, a função tira uma ``foto'' do estado do contêiner a cada 0,25 s
enquanto a inserção ocorre, e resume os recursos utilizados ao final.

\subsubsection{Medição de espaço em disco}
\label{subsec:medicao-espaco}

Além do tempo, o espaço ocupado é uma métrica relevante, pois um índice que
acelera as consultas pode ter um custo significativo de armazenamento. As
funções do Código~\ref{lst:medir-espaco}, no Apêndice~\ref{chapter:apendice-espaco}, medem o tamanho dos dados e o dos
índices em cada banco.

Nos dois bancos, o espaço é medido em disco e expresso em mebibytes (MiB). No PostgreSQL, o tamanho dos dados
vem de \texttt{pg\_table\_size}, que considera a tabela e a área TOAST.
Já o tamanho dos índices é a soma dos tamanhos dos índices da tabela, sem o índice da chave
primária. Como todos os demais índices são removidos antes de cada medição, o
valor obtido corresponde apenas ao índice criado para a consulta em teste.

No MongoDB, o tamanho dos dados vem do campo \texttt{storageSize} do comando
\texttt{collStats}, que reflete o espaço efetivamente alocado em disco, já com a
compressão aplicada pelo \textit{WiredTiger}. Em seguida, mede-se também o
tamanho sem compressão. O tamanho dos
índices vem do campo \texttt{totalIndexSize}. O índice da chave primária
(\texttt{\_id}) é descontado, pois existe em todos os cenários e não representa
o custo do índice em teste. 

O espaço dos dados é medido uma vez após cada carga. O espaço ocupado pelos índices é reportado apenas nas consultas em que há um índice criado
para o teste (C1, C4, C5, C6 e C7). As consultas C2, C3 e C8 localizam o item pelo
identificador, que já é coberto pelo índice da chave primária nos dois bancos.


\subsection{Experimentos de Escalabilidade}

Para observar como cada operação se comporta quando o cenário cresce, foram
definidos três experimentos de escalabilidade, cada um variando um fator e
mantendo os demais fixos: o volume total do banco, a seletividade do filtro e o
tamanho do lote de escrita. Os
resultados de cada experimento são gravados em arquivos JSON para análise
posterior.

\subsubsection{Volume total do banco}

Este experimento mede como cada operação se comporta quando o banco inteiro
cresce. A pergunta é se ao buscar exatamente a mesma coisa o tempo piora e, em
caso afirmativo, se piora na mesma proporção do crescimento dos dados. O banco é
avaliado em três volumes, apresentados na Tabela~\ref{tab:volumes}. Para cada volume, os bancos são recriados do zero e carregados com dados gerados com a mesma semente, de modo que o conteúdo é reproduzível.

\begin{table}[H]
    \centering
    \renewcommand{\arraystretch}{1.1}
    \caption{Volumes de dados avaliados.}
    \label{tab:volumes}
    \footnotesize
    \begin{tabular}{|l|c|c|}
        \hline
        \textbf{Volume} & \textbf{Restaurantes} & \textbf{Itens} \\ \hline
        Pequeno & 50 & 5.000 \\ \hline
        Médio & 200 & 100.000 \\ \hline
        Grande & 1.000 & 1.000.000 \\ \hline
    \end{tabular}
    \vspace{0.4em}
    \legend{Fonte: Elaborado pela autora.}
\end{table}

As consultas e os
parâmetros são sempre os mesmos, o que muda de uma rodada para outra é apenas o
tamanho do banco. Com isso, a proporção de registros que satisfazem cada filtro
permanece aproximadamente constante. 

\subsubsection{Seletividade}

A seletividade indica a fração dos registros que satisfaz o filtro de uma
consulta. Este experimento mantém o volume do banco fixo (Grande) e varia o valor
filtrado, de modo que a mesma consulta retorne quantidades diferentes de
registros. A quantidade de linhas retornadas é
registrada em cada medição, pois a seletividade depende dos dados gerados.

Foram escolhidas as consultas de leitura cujo filtro possui um parâmetro que
controla o tamanho do resultado: C1, C4, C5 e C7, mostrado na Tabela~\ref{tab:seletividade}. 
As outras consultas não foram incluídas por não terem um parâmetro variável ou pelo desejo de visualizar apenas um valor muito seletivo, como C6, que retorna apenas os itens com desconto de 80\%.

\begin{table}[H]
    \centering
    \renewcommand{\arraystretch}{1.1}
    \caption{Parâmetros variados no experimento de seletividade.}
    \label{tab:seletividade}
    \footnotesize
    \begin{tabular}{|l|p{4cm}|p{7cm}|}
        \hline
        \textbf{Consulta} & \textbf{Parâmetro variado} & \textbf{Valores} \\ \hline
        C1 & Categoria do item & Sobremesa (cerca de 10\% dos itens), bebida (cerca de 20\%) e pizza (cerca de 35\%) \\ \hline
        C5 & Largura da faixa de preço & R\$ 5--12, R\$ 5--25, R\$ 5--50 e R\$ 5--70 \\ \hline
        C4 & Porte do restaurante & Pequeno, médio e grande \\ \hline
        C7 & Opção buscada & Picles extra (cerca de 0,09\% dos itens), Cheddar extra (cerca de 9,4\%) e Bacon extra (cerca de 20,5\%) \\ \hline
    \end{tabular}
    \vspace{0.4em}
    \legend{Fonte: Elaborado pela autora.}
\end{table}

O objetivo é observar como o tempo e os índices respondem ao aumento do número de registros
retornados.

\subsubsection{Lote de escrita}

Este experimento verifica se o custo da operação \textit{por item} diminui conforme o lote
cresce e a partir de que tamanho essa melhoria deixa de ocorrer. O volume do
banco é fixo (Grande), e os lotes variam entre 1, 10, 100, 1.000, 10.000, 50.000 e 100.000 itens.

Quatro operações são avaliadas, cada uma sobre $N$ itens: C2 (como INSERT e UPDATE de forma separada), C3 e C8. Nos experimentos C3 e C8, os identificadores dos itens do lote são sorteados
entre os itens que já existem no banco. No C2, são gerados identificadores novos,
pois o \textit{upsert} insere itens que ainda não existem. Em todos os casos, o
lote inteiro é enviado ao banco em uma única operação. No PostgreSQL, essa
operação é uma instrução SQL seguida de \textit{commit} e no MongoDB, é um
\texttt{bulk\_write} (C2) ou um \texttt{update\_many} com \texttt{\$in} (C3 e C8).

Para que as repetições executem o mesmo trabalho,
utiliza-se duas funções auxiliares:

\begin{itemize}
  \item \texttt{antes}: é o parâmetro opcional da função \texttt{medir\_tempo}
  (Seção~\ref{subsec:medicao-tempo}), executado antes de cada repetição e fora
  da cronometragem. Em cada consulta ele desempenha a seguinte função:
  \begin{itemize}
    \item C2 (inserção): os itens do lote são apagados e são gerados novos eventos.
    \item C2 (atualização): os itens são criados uma vez, antes das repetições, e a cada 
    repetição recebem eventos novos, com versão maior e conteúdo diferente, de modo que 
    a atualização condicional seja de fato executada.
    \item C3 e C8: os campos alterados voltam aos valores originais.
  \end{itemize}
  \item \texttt{depois}: é executado uma única vez, ao fim das dez repetições de
  cada tamanho de lote, e devolve o banco ao estado original, de modo que o
  lote seguinte parta das mesmas condições.
\end{itemize}

Além do tempo total do lote, calcula-se o tempo por item, dado pela razão entre o
tempo médio do lote e o número de itens $N$. É importante ressaltar que o tamanho do lote não trata
de vários clientes simultâneos e sim de um único cliente enviando vários itens de uma vez.

\section{Resultados Obtidos}
\label{sec:resultados}

Os resultados vêm de uma única rodada completa dos experimentos, executados nos dois SGBDs com os mesmos dados sintéticos. Vale relembrar, que todas as medições de 
consulta e de escrita, os contêineres são reiniciados e o cache do sistema operacional é descartado antes de cada repetição. 
Cada medição de consulta é repetida 10 vezes e a carga inicial de cada volume é repetida 5 vezes, recriando os bancos do zero. 

As convenções dos gráficos são: PostgreSQL em azul e MongoDB em laranja. Em consultas onde tem variação de volume, elas são avaliadas em 5 mil, 100 mil e 1 milhão de itens. Caso tiver outro tipo de variação, 
o volume é fixado em 1 milhão de itens.

A forma como a média é mostrada depende do gráfico. 
Nos gráficos de variação de volume e de
lote, cada cenário traz uma faixa do primeiro quartil (Q1) ao terceiro quartil (Q3), um fio do
mínimo ao máximo e a média das 10 repetições, representadas, respectivamente, por um traço horizontal na faixa e um
ponto sobre a linha. A faixa de Q1 a Q3 cobre a metade central das repetições, deixando de fora os
25\% menores e os 25\% maiores valores. Essa faixa ajuda a verificar se a média resume bem os tempos medidos, já que
se a média está dentro da faixa, nenhum valor extremo a está distorcendo. Além
disso, quanto mais estreita a faixa, mais parecidos são os tempos e melhor a
média os representa.


Nos gráficos de carga inicial, a altura da barra é a média das
5 rodadas e o fio vai do mínimo ao máximo entre elas, como são poucas rodadas,
não se usa a faixa de Q1 a Q3. Já os de espaço ocupado, o tamanho foi o mesmo
em todas as rodadas, então não há variação a representar. Nos demais gráficos, as 10 repetições variaram pouco, por isso foi escolhido apresentar apenas a média dos valores, já que desenhar a variação sobre cada
barra apenas poluiria a figura.

\subsection{Efeito do Volume e do Índice}

Esta subseção trata da variação do volume do banco. Cada gráfico mostra uma
consulta, com o volume no eixo horizontal e o
tempo no vertical (em escala logarítmica). As consultas foram agrupadas pelo 
comportamento observado, apresenta-se no texto uma consulta de cada grupo, e as
demais estão no Apêndice~\ref{chapter:apendice-volume}. Com 5 mil itens as operações levam
poucos milissegundos, e pequenas oscilações pesam mais, fazendo com que a variação chegue a
cerca de 20\% a 50\% da média. Portanto a avaliação do efeito do volume é mais confiável nos volumes maiores, em que a variação cai para cerca de 5\% a 10\% da média. 

Além do volume, cada consulta é medida com e sem índice, para avaliar quanto o
índice ajuda. Esse efeito é medido pelo ganho do índice, e duas situações exigiram
cuidados no experimento:
\begin{itemize}
    \item \textbf{Índices de id:} as buscas por identificador (C2, C3 e C8) já são
    cobertas por um índice criado automaticamente pelos bancos (a chave primária no
    PostgreSQL e o \texttt{\_id} no MongoDB). Portanto, para medir o
    quanto ele ajuda, foi criado um mecanismo que faz o banco ignorá-lo e varrer a
    tabela inteira.
    \item \textbf{Índice forçado:} em algumas consultas (C1 e C5), o PostgreSQL
    decide por conta própria não usar o índice criado e faz uma varredura. Nesses
    casos, o experimento também mede a consulta obrigando o uso do índice, apenas
    para estudar o seu efeito.
\end{itemize}


\subsubsection{Consultas por id}

Essa consulta localiza o item pelo identificador e devolve uma única
linha, o que a torna a mais seletiva do experimento.
Sem índice, o tempo cresce com o volume porque o banco varre a tabela inteira:
o PostgreSQL passa de 35,4 ms (5 mil) para 111,5 ms
(100 mil) e 729,5 ms (1 milhão), e o MongoDB, de 25,9 ms para 124,0 ms e 1.052,9 ms.
Com índice, o tempo fica quase constante em qualquer volume, pois o índice leva
direto ao item. Nesse cenário, o MongoDB é quase duas vezes mais rápido que o PostgreSQL, em um volume médio e grande.
A C2 e a C8 (Apêndice~\ref{chapter:apendice-volume}) 
apresentam o mesmo comportamento.

\begin{figure}[H]
    \centering
    \caption{Efeito do volume no tempo de C3.}
    \includegraphics[width=0.7\textwidth]{grafico1_escala_C3_id.png}
    \label{fig:escala-c3}
    \legend{Fonte: Elaborado pela autora.}
\end{figure}

\subsubsection{Consultas Seletivas}

Essa consulta é seletiva, mas não tanto quanto as buscas por identificador, pois
devolve algumas linhas e não apenas uma. Sem índice, o tempo cresce com o volume: o
PostgreSQL passa de 36,6 ms (5 mil) para 109,9 ms (100 mil) e 1.175,8 ms (1 milhão), e
o MongoDB, de 19,3 ms para 117,3 ms e 1.075,9 ms. 
Com índice, diferente das consultas por id, o tempo aumenta com o volume, porque a
consulta devolve mais linhas quando o banco é maior: o PostgreSQL passa de 24,8 ms
para 49,5 ms e 218,7 ms, e o MongoDB, de 18,3 ms para 58,1 ms e 371,1 ms. Mesmo assim,
o índice ajuda bastante, isso faz com que o PostgreSQL seja cerca de 1,7 vez mais
rápido que o MongoDB com índice em 1 milhão de itens. A C4 (Apêndice~\ref{chapter:apendice-volume})
apresenta comportamento semelhante.


\begin{figure}[H]
    \centering
    \caption{Efeito de volume no tempo de C6.}
    \includegraphics[width=0.7\textwidth]{grafico1_escala_C6.png}
    \label{fig:escala-c6}
    \legend{Fonte: Elaborado pela autora.}
\end{figure}

\subsubsection{Consultas pouco Seletivas}

Essa consulta não é seletiva, pois devolve uma grande parte da tabela (cerca de 33\%) como resposta.
Sem índice, o tempo cresce com o volume: o
PostgreSQL passa de 41,7 ms (5 mil) para 140,7 ms (100 mil) e 1.182,9 ms (1 milhão), e
o MongoDB, de 26,0 ms para 174,2 ms e 1.645,6 ms. Com índice, o tempo cresce do mesmo
modo e o índice não ajuda. No PostgreSQL, o ponto crítico é e em 1 milhão de itens, onde o próprio banco decide 
ignorar o índice. Ao forçar o uso, o tempo vai a 2.526,9 ms (cerca de 2,1 vezes mais lento). No MongoDB, o
índice quase não altera o tempo. O melhor cenário para essa consulta acaba sendo o PostgreSQL sem índice. 
A C5 e a C7 (Apêndice~\ref{chapter:apendice-volume}) apresentam o
mesmo comportamento.


\begin{figure}[H]
    \centering
    \caption{Efeito de volume no tempo de C1.}
    \includegraphics[width=0.7\textwidth]{grafico1_escala_C1.png}
    \label{fig:escala-c1}
    \legend{Fonte: Elaborado pela autora.}
\end{figure}

\subsubsection{Síntese do Ganho do Índice}

A Figura~\ref{fig:ganho-indice} resume o ganho do índice por consulta. Esse ganho é o tempo sem índice dividido pelo tempo com índice. 
Acima de 1× o índice melhora a consulta e abaixo de 1× piora. O padrão é a seletividade: quanto menos tuplas a consulta devolve, mais
o índice ajuda. Nas buscas por id (C2, C3 e C8) o ganho é de dezenas de vezes, nas
seletivas (C4 e C6), de unidades a dezenas, e nas pouco seletivas (C1, C5 e C7),
próximo de 1$\times$ ou menor. Nas buscas por id, o ganho do MongoDB é maior, nas
seletivas, o do PostgreSQL.

\begin{figure}[H]
    \centering
    \caption{Síntese do ganho proporcionado pelo índice em cada consulta.}
    \includegraphics[width=0.7\textwidth]{grafico2_ganho_indice.png}
    \label{fig:ganho-indice}
    \legend{Fonte: Elaborado pela autora.}
\end{figure}

\subsection{Seletividade}

Nesta subseção varia-se o filtro de cada
consulta, de modo que ela devolva mais ou menos tuplas. Foram avaliadas a C1 e a C4,
apresentadas a seguir, que representam dois padrões diferentes. A C5 e
a C7 repetem esses padrões e, por isso, estão no
Apêndice~\ref{chapter:apendice-seletividade}.

Na C1, o filtro devolve de 10\% a 35\% da tabela, conforme a categoria buscada
(Figura~\ref{fig:seletividade-c1}). Nessa faixa, o índice não ajuda. No
PostgreSQL, para sobremesa, o tempo vai de 838 ms sem
índice para 5.431 ms com índice, e o mesmo ocorre para bebida. Já para pizza, o PostgreSQL abandona o
índice e faz uma varredura da tabela. No MongoDB o índice continua sendo usado, mas quase não altera o tempo.

\begin{figure}[H]
    \centering
    \caption{Efeito da seletividade no tempo de C1.}
    \includegraphics[width=0.7\textwidth]{grafico3_seletividade_C1.png}
    \label{fig:seletividade-c1}
    \legend{Fonte: Elaborado pela autora.}
\end{figure}


Na C4, o filtro devolve menos de 0,5\% da tabela, conforme o
porte do restaurante (Figura~\ref{fig:seletividade-c4}). Com índice: no PostgreSQL, passa de 70 ms
(restaurante pequeno) para 235 ms (médio) e 840 ms (grande), e no MongoDB, de 103 ms
para 380 ms e 571 ms. O índice ajuda nos três portes e nos dois bancos, mas o ganho
diminui conforme o restaurante tem mais itens, pois cada linha devolvida é uma
leitura a mais. Sem índice, o tempo praticamente não muda com o porte do restaurante.

\begin{figure}[H]
    \centering
    \caption{Efeito da seletividade no tempo de C4.}
    \includegraphics[width=0.7\textwidth]{grafico3_seletividade_C4.png}
    \label{fig:seletividade-c4}
    \legend{Fonte: Elaborado pela autora.}
\end{figure}


\subsection{Lote}

Nesta subseção varia-se o tamanho do lote, isto é, quantos itens são enviados de uma vez 
em uma única operação de escrita. Apresenta-se no texto a C2 (inserção)
(Figura~\ref{fig:lote-c2-insercao}), e as outras consultas que seguem o mesmo padrão estão no
Apêndice~\ref{chapter:apendice-lote}. O custo por item diminui muito conforme o lote aumenta. Com 1 item, cada operação custa de 14 a 28 ms, que é o custo fixo
de uma escrita confirmada. Quando o lote
cresce, esse custo passa, do lote de 1 item para o de 100 mil,
de 28,3 para 0,15 ms no PostgreSQL e de 13,7 para 0,011 ms no MongoDB.


Esse ganho não é sempre contínuo. No PostgreSQL, o custo por item deixa de cair em
torno de 50 mil itens e no lote de 100 mil tende a ficar igual ou pior. 
Isso sugere que lotes muito grandes passam a custar em memória e o ganho de dividir o custo fixo se esgota. No
MongoDB a saturação é mais tardia e continua caindo até 100 mil. O MongoDB é mais
rápido por item na maioria dos tamanhos de lote, e a vantagem aumenta nos lotes grandes.

\begin{figure}[H]
    \centering
    \caption{Efeito do tamanho do lote em C2.}
    \includegraphics[width=0.7\textwidth]{grafico4_lote_C2_insercao.png}
    \label{fig:lote-c2-insercao}
    \legend{Fonte: Elaborado pela autora.}
\end{figure}

\subsection{Carga Inicial dos Dados}

Esta subseção apresenta as três medidas que descrevem o custo de carregar os dados 
nos SGBDs pela primeira vez.
A (Figura~\ref{fig:carga-tempo}) mostra quanto tempo dura o envio dos itens ao banco. 
A (Figura~\ref{fig:carga-memoria}) mostra quanta
memória o contêiner ocupa em média durante a carga.
A (Figura~\ref{fig:carga-cpu}) mostra
quanto do processador do contêiner o banco usa durante a carga.
Vale ressaltar que os resultados com 5 mil itens não são considerados para a análise de CPU, já que as operações
levam poucos milissegundos, e pequenas oscilações do sistema pesam muito no tempo
total, o que gera uma variação de 20\% a 50\% da média entre as repetições.


Nos dois bancos, o tempo e a memória aumentam com o volume, mas em ritmos
diferentes. De 100 mil para 1 milhão de itens, os dados crescem 10 vezes e o tempo
de carga acompanha esse crescimento: passa de 4,9 para 63,7 s no PostgreSQL e de
1,5 para 16,6 s no MongoDB. A memória cresce bem menos: de 219 para 318 MB no
PostgreSQL e de 346 para 593 MB no MongoDB. A CPU, ao contrário dos dois, não
cresce com o volume, pois os lotes chegam ao banco no mesmo ritmo e o que aumenta é
a duração da carga.

Na comparação entre os bancos, o MongoDB carrega os dados de 3 a 4 vezes mais
rápido que o PostgreSQL em todos os volumes, e também usa menos CPU (cerca de 32\%
contra 57\%). Uma possível explicação para essas duas diferenças é a verificação da
chave estrangeira a cada item inserido, que o PostgreSQL faz e o MongoDB não. Em
contrapartida, o PostgreSQL ocupa menos memória: com 1 milhão de itens, são 318 MB
contra 593 MB do MongoDB.


\begin{figure}[H]
    \centering
    \caption{Tempo, uso de memória e uso de CPU durante a carga inicial dos dados.}
    \begin{subfigure}[t]{0.48\textwidth}
        \centering
        \includegraphics[width=\textwidth]{grafico6_carga_tempo.png}
        \caption{Tempo.}
        \label{fig:carga-tempo}
    \end{subfigure}\hfill
    \begin{subfigure}[t]{0.48\textwidth}
        \centering
        \includegraphics[width=\textwidth]{grafico8_carga_memoria.png}
        \caption{Memória.}
        \label{fig:carga-memoria}
    \end{subfigure}

    \vspace{1em}
    \begin{subfigure}[t]{0.48\textwidth}
        \centering
        \includegraphics[width=\textwidth]{grafico7_carga_cpu.png}
        \caption{CPU.}
        \label{fig:carga-cpu}
    \end{subfigure}
    \legend{Fonte: Elaborado pela autora.}
\end{figure}

\subsection{Armazenamento}

A Figura~\ref{fig:espaco-indice} mostra o tamanho dos índices criados para cada consulta. 
Os índices B-tree (C1, C4 e C5) ocupam pouco espaço, e os dos dois bancos têm tamanho parecido. 
Na C7, o índice GIN do PostgreSQL e o \textit{multikey} do MongoDB são praticamente iguais (17 e 18 MB), pois ambos guardam
uma entrada por elemento do array. A maior diferença está na C6: o GIN do PostgreSQL
ocupa 11 MB, e o índice \textit{wildcard} do MongoDB, 148 MB. Isso acontece, porque o GIN indexa apenas uma
coluna JSONB, enquanto o \textit{wildcard} indexa todos os campos de todos os
documentos. Isso o torna flexível, pois serve
para qualquer campo, mas bem mais caro, a ponto de ser maior que a própria coleção
comprimida em disco.

\begin{figure}[H]
    \centering
    \caption{Espaço em disco ocupado pelos índices.}
    \includegraphics[width=0.6\textwidth]{grafico5_espaco_indice.png}
    \label{fig:espaco-indice}
    \legend{Fonte: Elaborado pela autora.}
\end{figure}

A Figura~\ref{fig:carga-espaco} mostra o espaço ocupado pelos dados após a carga
inicial, que cresce proporcionalmente ao volume nos dois bancos. Sem considerar a
compressão, os dois ocupam praticamente o mesmo espaço. Porém, esse valor serve apenas
de comparação, pois o MongoDB comprime os dados por padrão e o PostgreSQL não. 
Logo, o que vale na prática é 108 MB no MongoDB contra 647 MB no PostgreSQL, o que é uma vantagem do MongoDB em uso real.

\begin{figure}[H]
    \centering
    \caption{Espaço em disco ocupado após a carga inicial dos dados.}
    \includegraphics[width=0.6\textwidth]{grafico9_carga_espaco.png}
    \label{fig:carga-espaco}
    \legend{Fonte: Elaborado pela autora.}
\end{figure}


\section{Dificuldades e Limitações}
\label{sec:dificuldades}

As principais dificuldades foram de execução e de comparabilidade. O notebook Jupyter, 
usado para desenvolver e explicar o experimento, não pôde ser
usado nos cenários de escalabilidade, pois o computador disponível travava durante a
execução. Por isso, a execução passou a ser feita por um \textit{script} Python equivalente,
no qual a rodada completa leva cerca de três horas.

Houve ainda uma mudança
de escopo, a proposta inicial envolvia o GraphQL, mas essa direção misturava as
camadas de acesso e de armazenamento, dificultando a visualização do que queria ser medido e analisado.
Por isso, a ideia foi focar nos SGBDs e em mecanismos que os caracterizam, como índices e escalabilidade.

Manter os dois bancos em condições comparáveis exigiu algumas decisões. A inserção é feita em
lotes de 10 mil itens, com \textit{commit} ao final de cada lote no PostgreSQL e
\texttt{journal=True} no MongoDB, de modo que as escritas tenham a mesma
durabilidade. Antes de cada repetição, os
contêineres são reiniciados e o cache do sistema operacional é descartado. O \textit{upsert} condicional (C2), não existe
como instrução única no MongoDB, foi implementado com uma lógica separada. 

Em relação as limitações, os experimentos foram executados em um único computador pessoal, com os bancos em contêineres de 2 CPUs e 3~GB cada, em um único nó e
com um único cliente, de forma sequencial. Por isso, os resultados podem não valer
para outro hardware ou para vários usuários simultâneos, e a capacidade da máquina
limitou o maior volume em 1 milhão de itens e o maior lote em 100 mil itens. Outra limitação é que os dados
são sintéticos, com semente fixa, e repetitivos, o que pode favorecer a compressão do
MongoDB. 

Os bancos usaram a configuração padrão, sem ajuste fino, e o tamanho do cache não é equivalente entre eles (128 MB no PostgreSQL contra cerca de 1 GB no MongoDB). Isso pode favorecer o MongoDB nas consultas que leem muitos dados de uma vez, embora o cache seja esvaziado antes de cada repetição.
Além disso, em 5 mil itens, a variação entre repetições é grande, fazendo com que a média não seja um valor confiável nesse volume.


\section{Considerações Finais}
\label{sec:consideracoes-finais-dev}

Este capítulo apresentou a metodologia e os resultados da comparação entre o
PostgreSQL com JSONB e o MongoDB, mostrando que o desempenho de cada um depende do
cenário e, em especial, da seletividade das consultas e do uso de índices. O
Capítulo~\ref{chapter:conclusao} apresenta as conclusões do trabalho e as
propostas de continuidade.
